\documentclass[11pt,answers]{exam} %noanswers
\usepackage{multicol}
\usepackage[usenames]{color} %used for font color
\usepackage{amssymb} %maths
\usepackage{amsmath} %maths
\usepackage[utf8]{inputenc} %useful to type directly diacritic characters
\usepackage{siunitx}
\usepackage{mathtools}
\usepackage{float}
\usepackage{systeme}
\usepackage[version=4]{mhchem}
\newcommand{\qtty}[2]{\left(\qty{#1}{#2}\right)}
\sisetup{display-per-mode = symbol }
\sisetup{inter-unit-product = \ensuremath { { } \cdot { } } }
\usepackage[letterpaper, total={7in, 10in}]{geometry}
\usepackage{graphicx} % Required for inserting images
\renewcommand{\epsilon}{\varepsilon}
\renewcommand{\vec}[1]{\mathbf{#1}}
\newcommand{\mat}[1]{\begin{bmatrix*}[r] #1 \end{bmatrix*}}
\newcommand{\dett}[1]{\begin{vmatrix*}[r] #1 \end{vmatrix*}}
\newcommand{\adj}[1]{\textrm{adj}\, #1}
\newcommand{\R}[1]{\mathbb{R}^{#1}}

\begin{document}
\flushleft
Name:\ \makebox[4in]{\hrulefill} \hspace{1in} Date:\ \makebox[1in]{\hrulefill}
\begin{center}
    \textbf{\S 4.1 Vector Spaces and Subspaces}
\end{center}

\begin{questions}
    \question Let $V$ be the first quadrant in the $xy-$plane; that is, let
    \begin{equation*}
        V= \left\{ \mat{x \\ y}:x\geq 0, y\geq 0\right\}
    \end{equation*}
    \begin{parts}
    \part If $\vec{u}$ and $\vec{v}$ are in $V$, is $\vec{u}+\vec{v}$ in $V$? Why?
    \part Find a specific vector $\vec{u}\in V$ and a specific sscalar $c$ such that $c\vec{u}\notin V$. (This is enough to show that $V$ is \emph{not} a vector space.)
    \end{parts}
    \question Let $W$ be the union of the first and third quadants in the $xy$-plane. That is, let
    \begin{equation*}
        W= \left\{ \mat{x \\ y}: xy\geq 0\right\}
    \end{equation*}
    \begin{parts}
    \part If $\vec{u}\in W$ and $c$ is any scalar, is $c\vec{u}$ in $W$? Why?
    \part Find specific vector $vec{u}$ and $\vec{v}$ in $A$ such that $\vec{u}+\vec{v}\notin W$. This is enough to that that $A$ is \emph{not} a vector space.

    \question Let $H$ be the set of points inside and on the unit circle in the $xy$-plane. That is, let 
    \end{parts}
    $
        \begin{aligned}[t]
            5x_1+7x_2 &= 3 \\
            2x_1 + 4x_2 &= 1
        \end{aligned}
    $
    \question
        $
        \begin{aligned}[t]
            4&x_1 + x_2 &= 6 \\
            3&x_1 + 2x_2 &= 5
        \end{aligned}
    $
        \question
        $
        \begin{aligned}[t]
            3x_1 - 2x_2 & 3 \\
            -4 x_1 + 6x_2 &= -5
        \end{aligned}
    $
        \question
        $
        \begin{aligned}[t]
            -5x_1 + 2x_2 &=9 \\
            3x_1 - x_2 &= -4
        \end{aligned}
    $
        \question
        $
        \begin{aligned}[t]
            x_1 + x_2 &= 3 \\
            -3x_1 + 2x_3 &= 0 \\
            x_2 - 2x_3 &= 2 
        \end{aligned}
    $
        \question
        $
        \begin{aligned}[t]
            x_1+3x_2+x_3&=8 \\
            -x_1 + 2x_3 &=4 \\
            3x_1 + x_2 &= 4
        \end{aligned}
    $

    \uplevel{In Exercsies 7-10, determine the values of the parameter $s$ for which the system has a unique solution and describe the solution.}
    \question
        $
        \begin{aligned}[t]
            6s &x_1 + 4 &x_2 &= 5 \\
            9 &x_1 + 2s &x_2 &= -2  
        \end{aligned}
        $
    \question 
        $
        \begin{alignedat}[t]{3}
            3s \, &x_1 + 5 &x_2 &= 3 \\
            12 &x_1 + 5s \, &x_2 &= 2
        \end{alignedat}
        $
    \question 
        $
        \begin{alignedat}[t]{3}
            s\, &x_1 + 2s \, &x_2 &= -1 \\
            3 &x_1 + 6s \, &x_2 &= 4
        \end{alignedat}
        $
    \question 
        $
        \begin{alignedat}[t]{3}
            s \, &x_1 - 2 &x_2 &= 1 \\
            4s \, &x_1 +4s\, &x_2 &= 2
        \end{alignedat}
        $
    
    \uplevel{In Exercises 11-16, compute the adjugate of the given matrix, and then use Theorem 8 to give the inverse of the matrix.}
    \question $\mat{0 & -2 & -1 \\ 5 & 0 & 0 \\ -1 & 1 & 1}$
    \question $\mat{1 & 1 & 3 \\ -2 & 2 & 1 \\ 0 & 1 & 1}$
    \question $\mat{3 & 5 & 4 \\ 1 & 0 & 1 \\ 2 & 1 & 1}$
    \question $\mat{1 & -1 & 2 \\ 0 & 2 & 1 \\ 3 & 0 & 6}$
    \question $\mat{5 & 0 & 0 \\ -1 & 1 & 0 \\ -2 & 3 & -1}$
    \question $\mat{1 & 2 & 4 \\ 0 & -3 & 1 \\ 0 & 0 & -2}$

    \question Show that if $A$ is $2 \times 2$, then Theorem 8 gives the same formula for $A^{-1}$ as that given by Theorem 4 in Seection 2.2.
    \question Suppose that all the entries in $A$ are integers and $\det{A}=1$. Explain why all the entries in $A^{-1}$ are integers.
    
    \uplevel{In Exercises 19-22, find the area of the parallelogram whose vertices are listed.}
    \question $(0,0), (5,2), (6,4), (11,6)$
    \question $(0,0), (-2,4),(6,-5),(4,-1)$
    \question $(-2,0), (0,3), (1,3), (-1,0)$
    \question $(0,-2), (5,-2), (-3,1), (2,1)$

    \question Find the volume of the parallelepiped with one vertex at the origin and adjacent vertices at $(1,0,-3), (1,2,4)$, and $(5,1,0)$.
    \question Find the volume of the parallelepiped with one vertex at the origin and adjacent vertices at $(1,3,0), (-2,0,2)$, and $(-1,3,-1)$.
    \question Use the concept of volume to explain why the determinant of a $3 \times 3$ matrix is zero if and only if $A$ is not invertible. Do not appeal to Theorem 4 in Section 3.2. Hint: Think about the columns of $A$.
    \question Let $T:\mathbb{R}^m \rightarrow \mathbb{R}^n$ be a linear transformation, and let $\vec{p}$ be a vector and $S$ a set in $\mathbb{R}^m$. Show that the image of $\vec{p}+S$ under $T$ is the translated set $T(\vec{p})+T(S)$ in $\mathbb{R}^n$.
    \question Let $S$ be the parallelogram determnined by the vectors $\vec{b}_1 = \mat{-2\\3}$ and $\vec{b}_2 = \mat{-2\\5}$, and let $A = \mat{6 & -3 \\ -3 & 2}$. Compute the area of the image of $S$ under the mapping $\vec{x} \mapsto A\vec{x}$.
    \question Let $S$ be the parallelogram determnined by the vectors $\vec{b}_1 = \mat{4 \\ -7}$ and $\vec{b}_2 = \mat{0 \\ 1}$, and let $A = \mat{5 & 2 \\ 1 & 1}$. Compute the area of the image of $S$ under the mapping $\vec{x} \mapsto A\vec{x}$.
    \question Find a formula for the area of a triangle whose vertices are $\vec{0}$, $\vec{v}_1$, and $\vec{v}_2$ in $\mathbb{R}^2$.
    \question Let $R$ be the triangle with vertices at $(x_1,y_1)$, $(x_2,y_2)$, and $(x_3,y_3)$. Show that 
    \begin{equation*}
        \left\{\textrm{area of triangle} \right\} = \frac{1}{2} \det{\mat{x_1 & y_1 & 1 \\ x_2 & y_2 & 1 \\ x_3 & y_3 & 1}}
    \end{equation*}
    Hint: Translate $R$ to the origin by subtracting one of the vertices, and use Exercise 29.

    \question Let $T:\mathbb{R}^3 \rightarrow \mathbb{R}^3$ be the linear transformation determined by the matrix $A = \mat{a & 0 & 0 \\ 0 & b & 0 \\ 0 & 0 & c}$, where $a$, $b$, and $c$ are positive numbers. Let $S$ be the unit ball, whose bounding surface has the equation $x_1^2 + x_2^2 + x_3^2 = 1$.
    \begin{parts}
        \part Show that $T(S)$ is bounded by the ellipsoid with the equation \begin{equation*} \frac{x_1^2}{a^2} + \frac{x_2^2}{b^2} + \frac{x_3^2}{c^2} = 1 \end{equation*}
        \part Use the fact that  the volume of the unit ball is $4\pi/3$ to determine the volume of the region bounded by the ellipsoid in part (a).
    \end{parts}
    \question Let $S$ be the tetrahedron in $\mathbb{R}^3$ with vertices at the vectors $\vec{0}$, $\vec{e}_1$, $\vec{e}_2$, and $\vec{e}_3$, and let $S^\prime$ be the tetrahedron with vertices at vectors $\vec{0}$, $\vec{v}_1$, $\vec{v}_2$, and $\vec{v}_3$. See the figure.
    \begin{parts}
        \part Describe a linear transformation that maps $S$ onto $S^\prime$.
        \part Find a formula for the volume of the tetrahedron $S^\prime$ using the fact that
        \begin{equation*}
            \left\{ \textrm{volume of $S$} \right\} = \frac{1}{3} \cdot \left\{ \textrm{area of base} \right\} \cdot \left\{ \textrm{height} \right\}
        \end{equation*}
    \end{parts}

    \question Let $A$ be an $n \times n$ matrix. If $A^{-1} = \frac{1}{\det{A}} \, \adj{A}$ is computed, what should $AA^{-1}$ be equal to in order to confirm that $A^{-1}$ has been found correctly?
    \question If a parallelogram fits inside a circle radius 1 and $\det{A} = 2$, where $A$ is a matrix whose columns correspond to the edges of the parallelogram, does it seem like $A$ and its determinant have been calculated correctly to correspond to the area of this parallelogram? Explain why or why not.

    \end{questions}

\end{document}